% === POEMA 1 ===
\begin{verse}
Qui, tenant l'ombre d'une brosse, \\
Nettoyait l'ombre d'un carrosse.
\end{verse}

% === POEMA 2 ===
\begin{verse}
'Cent défauts qu'on impute au siècle où tu naquis \\
Ne profaneroient pas tes ouvrages exquis\pagenote{'Exquis' is good.}.'
\end{verse}

% === POEMA 3 ===
\begin{verse}
'Clair de ruisseaux et sombre de verdure \\
Saisissait les esprits d'une secrete horreur; \\
\vin La simple et naive nature \\
\vin S'y faisoit voir si belle et si pure, \\
\vin Que mille fois il benit son erreur.'
\end{verse}

% === POEMA 4 ===
\begin{verse}
\vin 'Rempli de douces reveries \\
Qu'inspirent les grands bois, les eaux et les prairies.'
\end{verse}

% === POEMA 5 ===
\begin{verse}
'Cent et Cent fois ma Gouvernante \\
\vin Au lieu de Fables d'animaux\pagenote{The Fables d'animaux are probably even older than contes Diamonds and Toads. A Mouse and a Frog, as well as the Old Woman who survives as La Fée, take part in the tale as the Kaffirs tell it in The Story of Five Heads, in Theal's Kaffir Folk Lore, pp. 48, 49. The Kaffir story slides into a form of Beauty and the Beast. By some unexplained accident a story of Mlle. L'Heritier's L'Adroite Princesse slipped into editions of Perrault's Contes, in 1721, if not earlier, and holds its place even now.} \\
\vin M'a raconté les traits moraux \\
De cette Histoire surprenante.'
\end{verse}

% === POEMA 6 ===
\begin{verse}
'Quelquefois, cependant, c'est ce petit Marmot \\
Qui fera le bonheur de toute la famille!'
\end{verse}

% === POEMA 7 ===
\begin{verse}
'Mais le cher motif de leur joye, \\
Comme un conte de la Mère Oye, \\
Se trouvant fabuleux et faux, \\
Ils deviendront tous bien pénauts.'
\end{verse}

% === POEMA 8 ===
\begin{verse}
'Cent fois ma nourrice ou ma mie \\
M'ont fait ce beau recit le soir pres des tisons, \\
Je n'y fais qu'ajouter un peu de broderie.'
\end{verse}

% === POEMA 9 ===
\begin{verse}
'And sum thurch the bodi hadde wounde \\
Wives ther lay on childe bedde \\
Sum dede and sum awedde.'
\end{verse}

% === POEMA 10 ===
\begin{verse}
\vin\vin Nec conjugis umquam \\
Busta meae videam: neu sim tumulandus ab illa.
\end{verse}

% === POEMA 11 ===
\begin{verse}
'Au finement de cest escrit \\
K'én Romanz ai turné et dit, \\
Me numerai par remembraunce \\
Marie ai num, si sui de Fraunce.

* * * * *

Pur amur le cumte Willaume \\
Le plus vaillant de cest Royaume, \\
M'entremis de cest livre feire \\
E de l'Angleiz en Roman treire, \\
Ysopet apeluns ce livre \\
Qu'il traveilla et fist escrire; \\
De Griu en Latin le turna. \\
Li Rois Henris qui moult l'ama \\
Le translata puis en Engleiz \\
E jeo l'ai rimé en Franceiz.'
\end{verse}

% === POEMA 12 ===
\begin{verse}
'tel bec came li plereit \\
E cum li Huite cox aveit.'
\end{verse}

% === POEMA 13 ===
\begin{verse}
'Deus Oremanz unt ja perduz \\
Que nus n'en est a bien venuz,'
\end{verse}

% === POEMA 14 ===
\begin{verse}
'qui trop creient autrui parole.'
\end{verse}

% === POEMA 15 ===
\begin{verse}
'C'est un trésor qui n'embarrasse point.'
\end{verse}

% === POEMA 16 ===
\begin{verse}
'Would I could wish my wishes all to rest, \\
And know to wish the wish that should be best,
\end{verse}

% === POEMA 17 ===
\begin{verse}
'Soft on the fell \\
A shield-may sleepeth, \\
The lime-trees' red plague \\
Playing about her. \\
The sleep-thorn set Odin \\
Into that maiden \\
For her choosing in war \\
The one he willed not.'
\end{verse}

% === POEMA 18 ===
\begin{verse}
\vin 'Qui suivent les jeunes Demoiselles \\
Jusques dans les maisons, jusques dans les Ruelles!'
\end{verse}

% === POEMA 19 ===
\begin{verse}
'By cam' a cattie and ate it a' up my loesome, \\
\vin Loesome Lady! \\
And sae will I you—worrie, worrie, gnash, gnash, \\
\vin Said she, said she!'
\end{verse}

% === POEMA 20 ===
\begin{verse}
'Aux jeunes gens pour l'ordinaire, \\
L'industrie et le savoir-faire \\
Valent mieux que des biens acquis.'
\end{verse}

% === POEMA 21 ===
\begin{verse}
De riche appauvri Dieu te gard' \\
Et de croquant passé richard!
\end{verse}

% === POEMA 22 ===
\begin{verse}
\vin Tak me up, bane by bane \\
\vin And pit me aneth yon grey stane,
\end{verse}

% === POEMA 23 ===
\begin{verse}
'Twas late in the night and the bairns grat, \\
The Mother below the mouls heard that,—
\end{verse}

% === POEMA 24 ===
\begin{verse}
Fee, fie, fo, fum, \\
I smell the blood of some earthly one!\pagenote{Compare L'Oiseau Vert. Cosquin, Contes de Lorraine, i. 103.}
\end{verse}

% === POEMA 25 ===
\begin{verse}
Pouvois-je mieux choisir pour rendre vrai-semblable \\
\vin Ce que la Fable à d'incroyable? \\
\vin Et jamais Fée, au tems jadis \\
\vin Fit-elle à jeune Créature, \\
\vin Plus de dons, \& de dons exquis, \\
\vin Que vous en a fait la Nature.
\end{verse}

% === POEMA 26 ===
\begin{verse}
Attendre quelque temps pour avoir un Epoux, \\
\vin Riche bien-fait, galant \& doux, \\
\vin La chose est assez naturelle, \\
Mais l'attendre cent ans et toûjours en dormant, \\
\vin On ne trouve plus de femelle, \\
\vin Qui dormist si tranquillement. \\
La Fable semble encor vouloir nous faire entendre, \\
Que souvent de l'Hymen les agreables nœuds, \\
Pour estre differez n'en sont pas moins heureux, \\
\vin Et qu'on ne perd rien pour attendre; \\
\vin Mais le sexe avec tant d'ardeur, \\
\vin Aspire à la foy conjugale, \\
\vin Que je n'ay pas la force ny le cœur, \\
\vin De luy prescher cette Morale.
\end{verse}

% === POEMA 27 ===
\begin{verse}
On voit icy que de jeunes enfans, \\
\vin Sur tout de jeunes filles, \\
\vin Belles, bien faites, \& gentilles, \\
Font tres-mal d'écouter toute sorte de gens, \\
\vin Et que ce n'est pas chose étrange, \\
\vin S'il en est tant que le loup mange. \\
\vin Je dis le loup, car tous les loups; \\
\vin Ne sont pas de la mesme sorte; \\
\vin Il en est d'une humeur accorte, \\
\vin Sans bruit, sans fiel \& sans couroux, \\
\vin Qui privez, complaisans \& doux, \\
\vin Suivant les jeunes Demoiselles, \\
Jusque dans les maisons, jusque dans les ruelles; \\
Mais helas! qui ne sçait que ces Loups doucereux, \\
\vin De tous les Loups sont les plus dangereux.
\end{verse}

% === POEMA 28 ===
\begin{verse}
La curiosité malgré tous ses attraits, \\
\vin\vin Couste souvent bien des regrets; \\
On en voit tous les jours mille exemples paroistre, \\
C'est, n'en déplaise au sexe, un plaisir bien leger, \\
\vin Dés qu'on le prend il cesse d'estre, \\
\vin Et toûjours il couste trop cher.
\end{verse}

% === POEMA 29 ===
\begin{verse}
\vin Pour peu qu'on ait l'esprit sensé, \\
\vin Et que du Monde on sçache le grimoire, \\
\vin\vin On voit bien-tost que cette histoire \\
\vin\vin Est un conte du temps passé; \\
\vin Il n'est plus d'Epoux si terrible, \\
\vin Ny qui demande l'impossible; \\
\vin Fut-il mal-content \& jaloux, \\
\vin Prés de sa femme on le voit filer doux; \\
Et de quelque couleur que sa barbe puisse estre, \\
On a peine à juger qui des deux est le maistre.
\end{verse}

% === POEMA 30 ===
\begin{verse}
\vin Quelque grand que soit l'avantage, \\
\vin De joüir d un riche heritage \\
Venant à nous de pere en fils, \\
Aux jeunes gens pour l'ordinaire, \\
L'industrie \& le sçavoir faire, \\
Vallent mieux que des biens acquis.
\end{verse}

% === POEMA 31 ===
\begin{verse}
Si le fils d'un Meûnier, avec tant de vitesse, \\
\vin Gagne le cœur d'une Princesse, \\
Et s'en fait regarder avec des yeux mourans, \\
\vin C'est que l'habit, la mine \& la jeunesse, \\
\vin Pour inspirer de la tendresse, \\
N'en sont pas des moyens toûjours indifférens.
\end{verse}

% === POEMA 32 ===
\begin{verse}
\vin Les Diamans \& les Pistoles, \\
\vin\vin Peuvent beaucoup sur les Esprits; \\
\vin Cependant les douces paroles \\
Ont encor plus de force, \& sont d'un plus grand prix.
\end{verse}

% === POEMA 33 ===
\begin{verse}
\vin L'honnesteté couste des soins, \\
\vin Et veut un peu de complaisance, \\
\vin Mais tost ou tard elle a sa récompense, \\
Et souvent dans le temps qu'on y pense le moins.
\end{verse}

% === POEMA 34 ===
\begin{verse}
La beauté pour le sexe est un rare tresor, \\
\vin De l'admirer jamais on ne se lasse; \\
\vin Mais ce qu'on nomme bonne grace, \\
\vin Est sans prix, \& vaut mieux encor.

\vin C'est ce qu'à Cendrillon fit avoir sa Maraine, \\
\vin En la dressant, en l'instruisant, \\
\vin Tant \& si bien qu'elle en fit une Reine: \\
(Car ainsi sur ce Conte on va moralisant.)

\vin Belles, ce don vaut mieux que d'estre bien coëffées, \\
Pour engager un cœur, pour en venir à bout, \\
\vin La bonne grace est le vrai don des Fées; \\
Sans elle on ne peut rien, avec elle on peut tout.
\end{verse}

% === POEMA 35 ===
\begin{verse}
\vin\vin C'est sans doute un grand avantage, \\
\vin\vin D'avoir de l'esprit, du courage, \\
\vin\vin De la naissance, du bon sens, \\
\vin\vin Et d'autres semblables talens, \\
\vin\vin Qu'on reçoit du Ciel en partage; \\
\vin\vin Mais vous aurez beau les avoir, \\
Pour vostre avancement ce seront choses vaines; \\
\vin Si vous n'avez, pour les faire valoir, \\
\vin Ou des parrains ou des maraines.
\end{verse}

% === POEMA 36 ===
\begin{verse}
Ce que l'on voit dans cet écrit, \\
Est moins un conte en l'air que la verité même; \\
Tout est beau dans ce que l'on aime, \\
Tout ce qu'on aime a de l'esprit.
\end{verse}

% === POEMA 37 ===
\begin{verse}
\vin\vin Dans un objet où la Nature, \\
\vin\vin Aura mis de beaux traits, \& la vive peinture \\
\vin D'un teint où jamais l'Art ne sçauroit arriver, \\
Tous ces dons pourront moins pour rendre un cœur sensible, \\
\vin\vin Qu'un seul agrément invisible, \\
\vin\vin Que l'Amour y fera trouver.
\end{verse}

% === POEMA 38 ===
\begin{verse}
On ne s'afflige point d'avoir beaucoup d'enfans, \\
\vin Quand ils sont tous beaux, bien-faits \& bien grands, \\
\vin Et d'un exterieur qui brille; \\
\vin Mais si l'un d'eux est foible ou ne dit mot, \\
\vin On le méprise, on le raille, on le pille, \\
Quelquefois cependant c'est ce petit marmot \\
Qui fera le bonheur de toute la famille.
\end{verse}

% === POEMA 39 ===
\begin{verse}
Le conte de Peau d'Ane est ici raconté \\
\vin Avec tant de naïveté, \\
\vin Qu'il ne m'a pas moins divertie \\
Que quand, auprès du feu, ma nourrice ou ma mie \\
Tenoient en le faisant mon esprit enchanté. \\
On y voit par endroits quelques traits de satire, \\
\vin Mais qui, sans fiel et sans malignité, \\
A tous également font du plaisir à lire. \\
Ce qui me plaît encor dans sa simple douceur \\
\vin C'est qu'il divertit et fait rire, \\
\vin Sans que mère, époux, confesseur, \\
\vin Y puissent trouver à redire.
\end{verse}

% === POEMA 40 ===
\begin{verse}
Il est des gens de qui l'esprit guindé, \\
\vin Sous un front jamais deridé \\
\vin Ne souffre, n'approuve \& n'estime \\
\vin Que le pompeux \& le sublime; \\
\vin Pour moi, j'ose poser en fait \\
Qu'en de certains momens l'esprit le plus parfait \\
Peut aimer sans rougir jusqu'aux Marionnettes; \\
\vin Et qu'il est des tems \& des lieux \\
\vin Où le grave \& le serieux \\
\vin Ne vallent pas d'agreables sornettes. \\
\vin Pourquoi faut-il s'émerveiller \\
\vin Que la Raison la mieux sensée, \\
\vin Lasse souvent de trop veiller, \\
\vin Par des contes d'Ogre\pagenote{Homme Sauvage qui mangeoit les petits enfans.} \& de Fée \\
\vin Ingenieusement bercée, \\
\vin Prenne plaisir à sommeiller, \\
\vin Sans craindre donc qu'on me condamne \\
\vin De mal employer mon loisir, \\
Je vais, pour contenter vôtre juste desir, \\
Vous conter tout au long l'histoire de Peau D'Asne.
\end{verse}

% === POEMA 41 ===
\begin{verse}
\vin Il étoit une fois un Roi \\
\vin\vin Le plus grand qui fût sur la Terre, \\
\vin Aimable en Paix, terrible en Guerre, \\
\vin Seul enfin comparable à soi: \\
Ses voisins le craignoient, ses Etats étoient calmes, \\
\vin Et l'on voyoit de toutes parts \\
\vin Fleurir, à l'ombre de ses palmes \\
\vin Et les Vertus \& les beaux Arts. \\
Son aimable Moitié, sa Compagne fidelle, \\
\vin Etoit si charmante \& si belle, \\
Avoit l'esprit si commode \& si doux \\
\vin Qu'il étoit encor avec elle \\
\vin Moins heureux Roi qu'heureux espoux. \\
\vin De leur tendre \& chaste Hymenée, \\
\vin Plein de douceur \& d'agrement \\
Avec tant de vertus une fille étoit née, \\
\vin Qu'ils se consoloient aisement \\
\vin De n'avoir pas de plus ample lignée.

\vin Dans son vaste \& riche Palais, \\
\vin Ce n'étoit que magnificence, \\
Partout y fourmilloit une vive abondance \\
\vin De Courtisans \& de Valets; \\
\vin Il avoit dans son Escurie \\
Grands \& petits chevaux de toutes les façons, \\
\vin Couverts de beaux caparaçons \\
\vin Roides d'or \& de broderie; \\
Mais ce qui surprenoit tout le monde en entrant \\
\vin C'est qu'au lieu plus apparent, \\
Un maître Asne étailloit ses deux grandes oreilles, \\
\vin Cette injustice vous surprend, \\
Mais, lorsque vous sçaurez ses vertus nompareilles, \\
Vous ne trouverez pas que l'honneur fût trop grand. \\
\vin Tel et si net le forma la Nature \\
\vin Qu'il ne faisoit jamais d'ordure, \\
\vin Mais bien beaux Ecus au soleil \\
\vin Et Loüis de toute maniere \\
Qu'on alloit recuëillir sur la blonde litiere \\
\vin Tous les matins à son reveil.

\vin Or le Ciel qui par fois se lasse \\
\vin De rendre les hommes contents, \\
Qui toûjours à ses biens mêle quelque disgrace \\
\vin Ainsi que la pluye au beau tems, \\
\vin Permit qu'une aspre maladie \\
Tout à coup de la Reine attaquât les beaux jours. \\
\vin Par tout on cherche du secours, \\
Mais ni la Faculté qui le Grec étudie, \\
\vin Ni les Charlatans ayant cours, \\
Ne pûrent tous ensemble arrêter l'incendie \\
Que la fievre allumoit en s'augmentant toûjours.

\vin Arrivée a sa derniere heure, \\
\vin Elle dit au Roi son époux \\
\vin Trouvez bon qu'avant que je meure, \\
\vin J'exige une chose de vous \\
\vin C'est que s'il vous prenoit envie \\
De vous remarier quand je n'y serai plus... \\
\vin —Ha! dit le Roi, ces soins sont superflus, \\
\vin Je n'y songerai de ma vie, \\
\vin Soyez en repos là dessus. \\
\vin Je le croi bien, reprit la Reine, \\
Si j'en prens à témoin vôtre amour vehement, \\
\vin Mais pour m'en rendre plus certaine \\
\vin Je veux avoir vôtre serment, \\
Adouci toute fois par ce temperamment \\
Que si vous rencontrez une femme plus belle, \\
\vin Mieux faite \& plus sage que moi, \\
Vous pourrez franchement lui donner vôtre foi \\
\vin Et vous marier avec elle: \\
\vin Sa confiance en ses attraits \\
Lui faisoit regarder une telle promesse \\
\vin Comme un serment surpris avec adresse \\
\vin De ne se marier jamais. \\
Le Prince jura donc, les yeux baignez de larmes \\
\vin Tout ce que la Reine voulut; \\
\vin La Reine entre ses bras mourut, \\
Et jamais un Mari ne fit tant de vacarmes. \\
A l'ouïr sanglotter \& les nuits \& les jours, \\
On jugea que son deüil ne lui durerait guerre \\
\vin Et qu'il pleuroit ses defuntes Amours \\
Comme un homme pressé qui veut sortir d'affaire. \\
On ne se trompa point. Au bout de quelques mois \\
Il voulut proceder à faire un nouveau choix; \\
\vin Mais ce n'étoit pas chose aisée, \\
\vin Il falloit garder son serment \\
\vin Et que la nouvelle Epousée \\
\vin Eût plus d'attraits \& d'agrement \\
Que celle qu'on venoit de mettre au monument.

\vin Ni la Cour en beautez fertile, \\
\vin Ni la Campagne, ni la Ville, \\
\vin Ni les Royaumes d'alentour \\
\vin Dont on alla faire le tour, \\
\vin N'en pûrent fournir une telle, \\
\vin L'Infante seule étoit plus belle \\
\vin Et possedoit certains tendres appas \\
\vin Que la deffunte n'avoit pas. \\
\vin Le Roy le remarqua lui-même \\
\vin Et brûlant d'un amour extreme \\
\vin Alla follement s'aviser \\
Que par cette raison il devoit l'épouser. \\
\vin Il trouva même un Casuiste \\
Qui jugea que le cas se pouvoit proposer; \\
\vin Mais la jeune Princesse triste \\
\vin D'oüïr parler d'un tel amour, \\
\vin Se lamentoit \& pleuroit nuit \& jour.

\vin De mille chagrins l'ame pleine \\
\vin Elle alla trouver sa Maraine, \\
\vin Loin dans une grotte à l'écart \\
De Nacre \& de Corail richement étoffée; \\
\vin C'étoit une admirable Fée \\
\vin Qui n'eut jamais de pareille en son Art. \\
\vin Il n'est pas besoin qu'on vous die \\
Ce qu'étoit une Fée en ces bienheureux tems, \\
\vin Car je suis sûr que votre Mie \\
\vin Vous l'aura dit dez vos plus jeunes ans.

\vin Je sçay, dit-elle, en voyant la Princesse \\
\vin\vin Ce qui vous fait venir ici, \\
Je sais de votre cœur la profonde tristesse \\
\vin Mais avec moi n'ayez plus de souci. \\
\vin Il n'est rien qui vous puisse nuire \\
Pourvû qu'à mes conseils vous vous laissiez conduire, \\
Votre Pere, il est vrai, voudroit vous épouser; \\
\vin Ecouter sa folle demande \\
\vin Seroit une faute bien grande \\
Mais sans le contredire on le peut refuser.

\vin Dites-lui qu'il faut qu'il vous donne \\
\vin Pour rendre vos desirs contents, \\
Avant qu'à son amour vôtre cœur s'abandonne \\
Une Robe qui soit de la couleur du Tems. \\
Malgré tout son pouvoir et toute sa richesse, \\
Quoi que le Ciel en tout favorise ses vœux, \\
Il ne pourra jamais accomplir sa promesse.

\vin Aussi-tôt la jeune Princesse \\
L'alla dire en tremblant à son Pere amoureux \\
\vin Qui dans le moment fit entendre \\
\vin Aux Tailleurs les plus importans \\
Que s'ils ne lui faisoient, sans trop le faire attendre, \\
Une robe qui fût de la couleur du Temps, \\
Ils pouvoient s'assurer qu'il les feroit \\
\vin\vin Tous pendre.

\vin Le second jour ne luisoit pas encor \\
\vin Qu'on apporta la robe desirée; \\
\vin Le plus beau bleu de l'Empirée \\
N'est pas, lorsqu'il est ceint de gros nuages d'or \\
\vin D'une couleur plus azurée. \\
De joye \& de douleur l'Infante penetrée \\
\vin Ne sçait que dire ni comment \\
\vin Se derober à son engagement. \\
\vin Princesse demandez-en une, \\
\vin Lui dit sa Maraine tout bas, \\
\vin Qui plus brillante \& moins commune, \\
\vin Soit de la couleur de la Lune \\
\vin Il ne vous la donnera pas. \\
A peine la Princesse en eut fait la demande \\
\vin Que le Roi dit à son Brodeur, \\
Que l'astre de la Nuit n'ait pas plus de splendeur \\
Et que dans quatre jours sans faute on me la rende.

Le riche habillement fut fait au jour marqué \\
\vin Tel que le Roy s'en étoit expliqué \\
Dans les Cieux où la Nuit a deployé ses voiles, \\
La Lune est moins pompeuse en sa robe d'argent \\
Lors même qu'au milieu de son cours diligent \\
Sa plus vive clarté fait pâlir les étoiles.

La Princesse admirant ce merveilleux habit \\
Estoit à consentir presque deliberée, \\
\vin Mais, par sa Maraine inspirée \\
\vin Au Prince amoureux elle dit, \\
\vin Je ne sçaurois être contente \\
Que je n'aye une Robe encore plus brillante \\
\vin Et de la couleur du Soleil; \\
Le Prince qui l'aimoit d'un amour sans pareil \\
Fit venir aussi-tôt un riche Lapidaire \\
\vin Et lui commanda de la faire \\
D'un superbe tissu d'or \& de diamans, \\
Disant que s'il manquoit à le bien satisfaire, \\
Il le feroit mourir au milieu des tourmens. \\
Le Prince fut exempt de s'en donner la peine, \\
\vin Car l'ouvrier industrieux, \\
\vin Avant la fin de la semaine \\
\vin Fit apporter l'ouvrage precieux \\
\vin Si beau, si vif, si radieux \\
\vin Que le blond Amant de Climene \\
\vin Lorsque sur la voute des Cieux \\
\vin Dans son char d'or il se promene \\
D'un plus brillant éclat n'ébloüit pas les yeux.

L'Infante que ces dons achevent de confondre \\
A son Pere, à son Roi ne sçait plus que répondre; \\
Sa Maraine aussi-tôt la prenant par la main, \\
\vin Il ne faut pas, lui dit-elle à l'oreille, \\
\vin Demeurer en si beau chemin, \\
\vin Est-ce une si grande merveille \\
\vin Que tous ces dons que vous en recevez \\
\vin Tant qu'il aura l'Asne que vous sçavez \\
\vin Qui d'écus d'or sans cesse emplit sa bource; \\
Demandez-lui la peau de ce rare Animal, \\
\vin Comme il est toute sa resource, \\
Vous ne l'obtiendrez pas, ou je raisonne mal.

\vin Cette Fée étoit bien sçavante, \\
\vin Et cependant elle ignoroit encor \\
Que l'amour violent pourvû qu'on le contente, \\
\vin Conte pour rien l'argent \& l'or; \\
La peau fut galamment aussi tôt accordée \\
\vin Que l'Infante l'eut demandée.

\vin Cette Peau quand on l'apporta \\
\vin Terriblement l'epouvanta \\
Et la fit de son sort amerement se plaindre, \\
Sa Maraine survint \& lui representa \\
Que quand on fait le bien on ne doit jamais craindre; \\
\vin Qu'il faut laisser penser au Roy \\
\vin Qu'elle est tout à fait disposée \\
A subir avec lui la conjugale Loi; \\
Mais qu'au même moment seule \& bien deguisée \\
Il faut qu'elle s'en aille en quelque Etat lointain \\
Pour éviter un mal si proche \& si certain.

Voici, poursuivit-elle, une grande cassette \\
\vin Où nous mettrons tous vos habits \\
\vin Vôtre miroir, vôtre toillette, \\
\vin Vos diamans \& vos rubis. \\
\vin Je vous donne encor ma Baguette; \\
\vin En la tenant en vôtre main \\
La cassette suivra vôtre même chemin. \\
\vin Toujours sous la Terre cachée; \\
\vin Et lorsque vous voudrez l'ouvrir \\
A peine mon bâton la Terre aura touchée \\
Qu'aussi-tôt à vos yeux elle viendra s'offrir.

\vin Pour vous rendre méconnaissable \\
La dépoüille de l'Asne est un masque admirable \\
\vin Cachez-vous bien dans cette peau, \\
On ne croira jamais, tant elle est effroyable \\
\vin Qu'elle renferme rien de beau.

\vin La Princesse ainsi travestie \\
De chez la sage Fée à peine fut sortie, \\
\vin Pendant la fraîcheür du matin \\
\vin Que le Prince qui pour la Fête \\
\vin De son heureux Hymen s'apprête \\
Apprend tout effrayé son funeste destin. \\
Il n'est point de maison, de chemin, d'avenuë \\
\vin Qu'on ne parcoure promptement, \\
On ne peut deviner ce qu'elle est devenuë.

Par tout se répandit un triste \& noir chagrin \\
\vin Plus de Nopces, plus de Festin, \\
\vin Plus de Tarte, plus de Dragées, \\
Les Dames de la Cour toutes découragées \\
\vin N'en dînerent point la plûpart; \\
Mais du Curé sur tout la tristesse fut grande, \\
\vin Car il en dejeuna fort tard \\
\vin Et qui pis est n'eut point d'offrande.

L'Infante cependant poursuivoit son chemin \\
Le visage couvert d'une vilaine crasse \\
\vin A tous Passans elle tendoit la main \\
Et tâchoit pour servir de trouver une place; \\
Mais les moins delicats \& les plus malheureux \\
La voyant si maussade \& si pleine d'ordure \\
Ne vouloient écouter ni retirer chez eux \\
\vin Une si sale creature. \\
Elle alla donc bien loin, bien loin, encor plus loin, \\
Enfin elle arriva dans une Metairie \\
\vin Où la Fermiere avoit besoin \\
\vin D'une soüillon, dont l'industrie \\
Allât jusqu'à sçavoir bien laver des torchons \\
\vin Et nettoyer l'auge aux Cochons.

On la mit dans un coin au fond de la cuisine \\
\vin Où les Valets, insolente vermine, \\
\vin Ne faisoient que la tirailler, \\
\vin La contredire \& la railler, \\
\vin Ils ne sçavoient quelle piece lui faire \\
\vin La harcelant à tout propos; \\
\vin Elle étoit la butte ordinaire \\
De tous leurs quolibets \& de tous leurs bons mots.

Elle avoit le Dimanche un peu plus de repos, \\
Car ayant du matin fait sa petite affaire, \\
Elle entroit dans sa chambre \& tenant son huis clos, \\
Elle se decrassoit, puis ouvroit sa cassette, \\
\vin Mettoit proprement sa toilette \\
\vin Rangeoit dessus ses petits pots, \\
Devant son grand miroir contente \& satisfaite; \\
De la Lune tantôt, la robe elle mettoit \\
Tantôt celle où le feu du Soleil éclattoit, \\
\vin Tantôt la belle robe blüe \\
Que tout l'azur des Cieux ne sçauroit égaler, \\
Avec ce chagrin seul que leur traînante queüe \\
Sur le plancher trop court ne pouvoit s'étaler. \\
Elle aimoit à se voir jeune, vermeille \& blanche \\
Et plus brave cent fois que nulle autre n'êtoit; \\
\vin Ce doux plaisir la sustentoit \\
\vin Et la menoit jusqu'à l'autre Dimanche.

\vin J'oubliois à dire en passant \\
\vin Qu'en cette grande Metairie \\
\vin D'un Roy magnifique \& puissant \\
\vin Se faisoit la Menagerie, \\
\vin Que là, Poules de Barbarie, \\
\vin Rales, Pintades, Cormorans, \\
\vin Oisons musquez, Cannes Petieres \\
Et mille autres oiseaux de bijares manieres, \\
\vin Entre eux presque tous differents \\
Remplissoient à l'envi dix cours toutes entieres.

\vin Le fils du Roy dans ce charmant sejour \\
\vin Venoit souvent au retour de la Chasse \\
\vin Se reposer, boire à la glace \\
\vin Avec les Seigneurs de sa Cour. \\
\vin Tel ne fut point le beau Cephale; \\
Son air étoit Royal, sa mine martiale \\
Propre à faire trembler les plus fiers bataillons; \\
Peau d'Asne de fort loin le vit avec tendresse \\
\vin Et reconnut par cette hardiesse \\
\vin Que sous sa crasse \& ses haillons \\
Elle gardoit encor le cœur d'une Princesse.

\vin Qu'il a l'air grand, quoi qu'il l'ait negligé, \\
\vin Qu'il est aimable, disoit-elle, \\
\vin Et que bienheureuse est la belle \\
\vin A qui son cœur est engagé. \\
D'une robe de rien s'il m'avoit honorée. \\
\vin Je m'en trouverois plus parée \\
\vin Que de toutes celles que j'ai.

Un jour le jeune Prince errant à l'aventure \\
\vin De bassecour en bassecour, \\
\vin Passa dans une allée obscure \\
\vin Où de Peau d'Asne étoit l'humble sejour. \\
Par hasard il mit l'œil au trou de la serrure; \\
\vin Comme il étoit fête ce jour \\
\vin Elle avoit pris une riche parure \\
\vin Et ses superbes vêtemens \\
Qui tissus de fin or \& de gros diamans \\
Egaloient du Soleil la clarté la plus pure. \\
\vin Le Prince au gré de son désir \\
\vin La contemple \& ne peut qu'à peine, \\
\vin En la voyant, reprendre haleine, \\
\vin Tant il est comblé de plaisir. \\
Quels que soient les habits, la beauté du visage, \\
\vin Son beau tour, sa vive blancheur, \\
\vin Ses traits fins, sa jeune fraîcheur \\
\vin Le touchent cent fois davantage, \\
\vin Mais un certain air de grandeur \\
Plus encore une sage \& modeste pudeur \\
Des beautez de son ame, asseuré témoignage, \\
\vin S'emparerent de tout son cœur.

Trois fois dans la chaleur du feu qui le transporte \\
\vin Il voulut enfoncer la porte, \\
\vin Mais croyant voir une Divinité, \\
Trois fois par le respect son bras fut arrêté, \\
Dans le Palais pensif il se retire \\
\vin Et là nuit \& jour il soupire, \\
\vin Il ne veut plus aller au Bal \\
\vin Quoi qu'on soit dans le Carnaval, \\
\vin Il hait la Chasse, il hait la Comedie \\
Il n'a plus d'appetit, tout lui fait mal au cœur \\
\vin Et le fond de sa maladie \\
\vin Est une triste \& mortelle langueur.

Il s'enquit quelle étoit cette Nymphe admirable \\
\vin Qui demeuroit dans une bassecour \\
\vin Au fond d'une allée effroyable, \\
\vin Où l'on ne voit goutte en plein jour. \\
C'est, lui dit-on, Peau d'Asne, en rien Nymphe ni bele \\
\vin Et que Peau d'Asne l'on appelle, \\
A cause de la peau qu'elle met sur son cou; \\
\vin De l'Amour c'est le vrai remede, \\
\vin La bête en un mot la plus laide, \\
\vin Qu'on puisse voir aprés le Loup: \\
\vin On a beau dire, il ne sçauroit le croire, \\
\vin Les traits que l'amour a tracez \\
\vin Toujours presens à sa memoire \\
\vin N'en seront jamais effacez.

\vin Cependant la Reyne sa Mere, \\
Qui n'a que lui d'enfant pleure \& se desespere, \\
De declarer son mal elle le presse en vain, \\
\vin Il gemit, il pleure, il soupire, \\
\vin Il ne dit rien, si ce n'est qu'il desire \\
Que Peau d'Asne lui fasse un gâteau de sa main; \\
Et la Mere ne sçait ce que son Fils veut dire; \\
\vin O Ciel! Madame, lui dit-on, \\
\vin Cette Peau d'Asne est une noire Taupe \\
\vin Plus vilaine encore \& plus gaupe \\
\vin Que le plus sale Marmiton. \\
N'importe, dit la Reyne, il le faut satisfaire, \\
Et c'est à cela seul que nous devons songer; \\
Il auroit eu de l'or, tant l'aimoit cette Mere, \\
\vin S'il en avoit voulu manger.

\vin Peau d'Asne donc prend sa farine \\
\vin Qu'elle avoit fait blutter exprés, \\
\vin Pour rendre sa pâte plus fine, \\
\vin Son sel, son beurre \& ses œufs frais, \\
\vin Et pour bien faire sa galette \\
\vin S'enferme seule en sa chambrette.

\vin D'abord elle se decrassa \\
\vin Les mains, les bras \& le visage, \\
Et prit un corps d'argent que vîte elle laça \\
\vin Pour dignement faire l'ouvrage, \\
\vin Qu'aussi-tôt elle commença.

On dit qu'en travaillant un peu trop à la hâte, \\
De son doigt par hazard il tomba dans la pâte \\
\vin Un de ses anneaux de grand prix, \\
Mais ceux qu'on tient sçavoir le fin de cette histoire \\
Asseurent que par elle exprés il y fut mis; \\
Et pour moi franchement, je l'oserois bien croire, \\
Fort seur que quand le Prince à sa porte aborda \\
\vin Et par le trou la regarda, \\
\vin\vin Elle s'en étoit apperçûë. \\
\vin Sur ce point la Femme est si druë, \\
\vin Et son œil va si promptement \\
\vin Qu'on ne peut la voir un moment, \\
\vin Qu'elle ne sçache qu'on l'a veüe. \\
Je suis bien seur encore, et j'en ferois serment \\
Qu'elle ne douta point que de son jeune Amant \\
\vin La Bague ne fût bien receuë.

On ne pêtrit jamais un si friand morceau, \\
Et le Prince trouva la galette si bonne \\
Qu'il ne s'en fallut rien que d'une faim gloutonne \\
\vin Il n'avalât aussi l'anneau. \\
\vin Quand il en vit l'émeraude admirable, \\
\vin Et du jonc d'or le cercle étroit, \\
\vin Qui marquoit la forme du doigt, \\
Son cœur en fut touché d'une joye incroyable; \\
\vin Sous son chevet il le mit à l'instant \\
\vin Et son mal toujours augmentant \\
\vin Les Medecins sages d'experience, \\
\vin En le voyant maigrir de jour en jour \\
\vin Jugerent tous par leur grande science \\
\vin Qu'il étoit malade d'amour.

\vin Comme l'Hymen, quelque mal qu'on en die, \\
Est un remede exquis pour cette maladie, \\
\vin On conclut à le marier; \\
\vin Il s'en fit quelque tems prier, \\
Puis dit, je le veux bien, pourvû que l'on me donne \\
\vin En mariage la personne \\
\vin Pour qui cet anneau sera bon; \\
\vin A cette bijare demande \\
De la Reine \& du Roi la surprise fut grande, \\
Mais il étoit si mal qu'on n'osa dire non. \\
\vin Voilà donc qu'on se met en quête \\
De celle que l'anneau, sans nul égard du sang, \\
\vin Doit placer dans un si haut rang, \\
\vin Il n'en est point qui ne s'apprête \\
\vin A venir presenter son doigt \\
\vin Ni qui veüille ceder son droit.

Le bruit ayant couru que pour prétendre au Prince, \\
\vin Il faut avoir le doigt bien mince, \\
\vin Tout Charlatan, pour être bien venu, \\
Dit qu'il a le secret de le rendre menu, \\
\vin L'une en suivant son bizare caprice \\
\vin Comme une rave le ratisse, \\
\vin L'autre en couppe un petit morceau, \\
Une autre en le pressant croit qu'elle l'appetisse, \\
\vin Et l'autre avec de certaine eau \\
Pour le rendre moins gros en fait tomber la peau; \\
\vin Il n'est enfin point de manœuvre \\
\vin Qu'une Dame ne mette en œuvre, \\
Pour faire que son doigt quadre bien à l'anneau.

L'essai fut commencé par les jeunes Princesses \\
\vin Les Marquises \& les Duchesses, \\
\vin Mais leurs doigts quoi que delicats \\
\vin Estoient trop gros \& n'entroient pas. \\
\vin Les Comtesses \& les Baronnes, \\
\vin Et toutes les nobles Personnes, \\
Comme elles tour à tour presenterent leur main \\
\vin Et la presenterent en vain.

\vin Ensuite vinrent les Grisettes, \\
\vin Dont les jolis \& menus doigts, \\
\vin Car il en est de tres-bien faites, \\
Semblerent à l'anneau s'ajuster quelquefois \\
Mais la Bague toujours trop petite ou trop ronde \\
D'un dedain presque égal rebuttoit tout le monde.

\vin Il fallut en venir enfin \\
\vin Aux Servantes, aux Cuisinieres, \\
\vin Aux Tortillons, aux Dindonnieres; \\
\vin En un mot à tout le fretin, \\
\vin Dont les rouges \& noires pattes, \\
\vin Non moins que les mains delicates \\
\vin Esperoient un heureux destin. \\
\vin Il s'y presenta mainte fille \\
\vin Dont le doigt gros \& ramassé, \\
Dans la Bague du Prince eût aussi peu passé \\
\vin Qu'un cable au travers d'une aiguille. \\
\vin On crut enfin que c'étoit fait, \\
\vin Car il ne restoit en effet, \\
Que la pauvre Peau d'Asne au fond de la cuisine, \\
\vin Mais comment croire, disoit-on, \\
\vin Qu'à regner le Ciel la destine, \\
\vin Le Prince dit, \& pourquoi non? \\
Qu'on la fasse venir. Chacun se prît à rire \\
Criant tout haut que veut-on dire, \\
De faire entrer ici cette sale guenon \\
Mais lorsqu'elle tira de dessous sa peau noire \\
Une petite main qui sembloit de l'yvoire, \\
\vin Qu'un peu de pourpre a coloré, \\
\vin Et que de la bague fatale, \\
\vin D'une justesse sans égale \\
\vin Son petit doigt fut entouré, \\
\vin La Cour fut dans une surprise \\
\vin Qui ne peut pas être comprise.

On la menoit au Roi dans ce transport subit, \\
Mais elle demanda qu'avant que de paraître \\
\vin Devant son Seigneur \& son Maître \\
On lui donnât le temps de prendre un autre habit, \\
\vin De cet habit, pour la verité dire, \\
\vin De tous côtez on s'apprétoit à rire, \\
Mais lorsqu'elle arriva dans les Appartemens \\
\vin Et qu'elle eut traversé les salles \\
\vin Avec ses pompeux vêtemens \\
Dont les riches beautez n'eurent jamais d'égales, \\
\vin Que ses aimables cheveux blonds \\
Mêlez de diamans dont la vive lumiere \\
\vin En faisoit autant de rayons, \\
\vin Que ses yeux bleus, grands, doux \& longs, \\
\vin Qui pleins d'une Majesté fiere \\
Ne regardent jamais sans plaire \& sans blesser, \\
Et que sa taille enfin si menüe \& si fine \\
Qu'avecque ses deux mains on eût pu l'embrasser, \\
Montrerent leurs appas \& leur grace divine; \\
Des Dames de la Cour, \& de leurs ornemens \\
\vin Tomberent tous les agrémens.

Dans la joye \& le bruit de toute l'Assemblée, \\
\vin Le bon Roi ne se sentoit pas \\
\vin De voir sa Bru posseder tant d'appas, \\
\vin La Reyne en étoit affolée, \\
\vin Et le Prince son cher Amant, \\
\vin De cent plaisirs l'âme comblée \\
Succomboit sous le poids de son ravissement. \\
Pour l'Hymen aussitôt chacun prit ses mesures, \\
Le Monarque en pria tous les Rois d'alentour, \\
\vin Qui tous brillans de diverses parures \\
Quitterent leurs Etats pour être à ce-grand jour \\
On en vit arriver des climats de l'Aurore, \\
\vin Montez sur de grands Elephans, \\
\vin Il en vint du rivage More, \\
\vin Qui plus noirs \& plus laids encore, \\
\vin Faisoient peur aux petits enfans; \\
\vin Il en debarque \& la Cour en abonde.

\vin Mais nul Prince, nul Potentat, \\
\vin N'y parut avec tant d'éclat \\
\vin Que le Pere de l'Epousée, \\
\vin Qui d'elle autrefois amoureux \\
Avoit avec le temps purifié les feux \\
\vin Dont son ame étoit embrasée, \\
Il en avoit banni tout desir criminel \\
\vin Et de cette odieuse flamme \\
\vin Le peu qui restoit dans son ame \\
N'en rendoit que plus vif son amour paternel. \\
\vin Dés qu'il la vit, que benit soit le Ciel \\
\vin Qui veut bien que je te revoye, \\
Ma chere enfant, dit-il, \&, tout pleurant de joye \\
\vin Courut tendrement l'embrasser; \\
Chacun à son bonheur voulut s'interesser, \\
Et le futur Espoux étoit ravi d'apprendre \\
Que d'un Roi si puissant il devenoit le Gendre. \\
\vin Dans ce moment la Maraine arriva \\
\vin Qui raconta toute l'histoire, \\
\vin Et par son recit acheva \\
\vin De combler Peau d'Asne de gloire.

\vin Il n'est pas malaisé de voir \\
Que le but de ce conte est qu'un Enfant apprenne \\
Qu'il vaut mieux s'exposer à la plus rude peine \\
\vin Que de manquer à son devoir.

\vin Que la Vertu peut être infortunée \\
\vin Mais qu'elle est toujours couronnée.

Que contre un fol amour \& ses fougueux transports \\
La Raison la plus forte est une foible digue, \\
\vin Et qu'il n'est point de si riches thresors \\
\vin Dont un Amant ne soit prodigue.

\vin Que de l'eau claire \& du pain bis \\
\vin Suffisent pour la nourriture \\
\vin De toute jeune Creature, \\
\vin Pourvu qu'elle ait de beaux habits.

\vin Que sous le Ciel il n'est point de femelle \\
\vin Qui ne s'imagine être belle, \\
\vin Et qui souvent ne s'imagine encor \\
Que si des trois Beautez la fameuse querelle, \\
\vin S'étoit demêlée avec elle \\
\vin Elle auroit eu la pomme d'or. \\
Le Conte de Peau d'Asne est difficile à croire, \\
Mais tant que dans le Monde on aura des Enfans, \\
\vin Des Meres \& des Meres-grands, \\
\vin On en gardera la memoire.
\end{verse}

% === POEMA 42 ===
\begin{verse}
Si vous étiez moins raisonnable, \\
\vin Je me garderois bien de venir vous conter \\
\vin La folle \& peu galante Fable, \\
\vin Que je m'en vais vous debiter. \\
Une aune de Boudin en fournit la matiere. \\
\vin Une aune de Boudin, ma chere: \\
\vin Quelle pitié! c'est une horreur, \\
\vin S'écrieroit une Pretieuse, \\
\vin Qui toujours tendre \& serieuse, \\
Ne veut ouir parler que d'affaires de cœur.

\vin Mais vous, qui mieux qu'autre qui vive, \\
\vin Sçavez charmer en racontant, \\
Et dont l'expression est toujours si naïve, \\
\vin Que l'on croit voir ce qu'on entend, \\
\vin Qui sçavez que c'est la maniere \\
\vin Dont quelque chose est inventé, \\
\vin Qui beaucoup plus que la matiere, \\
\vin De tout recit fait la beauté, \\
Vous aimerez ma Fable \& sa moralité; \\
J'en ai, j'ose le dire, une assurance entiere.

Il étoit une fois un pauvre Bucheron, \\
\vin Qui las de sa penible vie, \\
\vin Avoit, disoit-il, grande envie \\
De s'aller reposer aux bords de l'Acheron, \\
\vin Representant dans sa douleur profonde, \\
\vin Que depuis qu'il étoit au monde, \\
\vin Le Ciel cruel n'avoit jamais \\
Voulu remplir un seul de ses souhaits.

Un jour que dans le bois il se mit à se plaindre, \\
A lui la foudre en main Jupitre s'apparut. \\
\vin On auroit peine à bien dépeindre \\
\vin La peur que le bonhomme en eut. \\
Je ne veux rien, dit-il, en se jettant par terre, \\
\vin Point de souhaits, point de Tonnerre, \\
\vin Seigneur, demeurons but à but. \\
\vin Cesse d'avoir aucune crainte, \\
Je viens, dit Jupiter, touché de ta complainte, \\
\vin Te faire voir le tort que tu me fais. \\
\vin Ecoute donc, je te promets, \\
Moi qui du monde entier suis le Souverain Maître, \\
D'exaucer pleinnement les trois premiers souhaits \\
Que tu voudras former sur quoi que ce puisse être \\
\vin Voi ce qui peut te rendre heureux, \\
\vin Voi ce qui peut te satisfaire, \\
Et comme ton bonheur dépend tout de tes vœux, \\
\vin Songes y bien avant que de les faire.

A ces mots Jupiter dans les Cieux remonta, \\
Et le gay Bucheron embrassant sa falourde, \\
Pour retourner chez lui sur son dos la jetta. \\
Cette charge jamais ne lui parut moins lourde, \\
\vin Il ne faut pas, disoit-il en trottant, \\
\vin Dans tout ceci rien faire à la legere \\
\vin Il faut, le cas est important, \\
En prendre avis de nôtre Menagere, \\
C'a, dit-il en entrant sous son toit de feugere, \\
\vin Faisons, Fanchon, grand feu, grand'chere, \\
\vin Nous sommes riches desormais, \\
\vin Et nous n'avons qu'à faire des souhaits. \\
Là dessus fort au long tout se fait il lui conte. \\
\vin A ce recit, l'Epouse vive \& prompte, \\
Forma dans son esprit mille vastes projets, \\
\vin Mais considerant l'importance \\
\vin De s'y conduire avec prudence, \\
Blaise, mon cher Ami, dit-elle à son Epoux, \\
\vin Ne gâtons rien par nôtre impatience, \\
\vin Examinons bien entre nous \\
\vin Ce qu'il faut faire en pareille occurrence. \\
Remettons à demain notre premier souhait, \\
\vin Et consultons nôtre chevet. \\
Je l'entens bien ainsi, dit le bonhomme Blaise, \\
Mais va tirer du vin derriere ces fagots. \\
A son retour il but, \&, goütant à son aise \\
\vin Pres d'un grand feu la douceur du repos, \\
Il dit, en s'appuyant sur le dos de sa chaise, \\
Pendant que nous avons une si bonne braise, \\
Une aune de Boudin viendroit bien à propos. \\
A peine acheva-t-il de prononcer ces mots, \\
Que sa Femme apperceut, grandement étonnée, \\
\vin Un Boudin fort long, qui partant \\
\vin D'un des coins de la cheminée, \\
\vin S'approchoit d'elle en serpentant. \\
\vin Elle fit un cri dans l'instant, \\
\vin Mais jugeant que cette avanture \\
\vin Avoit pour cause le souhait \\
\vin Que par bêtise toute pure \\
\vin Son homme imprudent avoit fait, \\
\vin Il n'est point de pouille, ni d'injure, \\
\vin Que de depit \& de couroux \\
\vin Elle ne dît a son Epoux.

Quand on peut, disoit-elle, obtenir un Empire, \\
\vin De l'or, des Perles, des Rubis, \\
\vin Des Diamans, de beaux Habits, \\
Est-ce alors du Boudin qu'il faut que l'on desire? \\
Eh bien, j'ai tort, dit-il, j'ai mal placé mon choix. \\
\vin J'ai commis une faute énorme, \\
\vin Je ferai mieux une autrefois. \\
Bon, bon, dit-elle, attendez-moi sous l'orme. \\
Pour faire un tel souhait, il faut être bien Bœuf. \\
L'Epoux plus d'une fois emporté de colere \\
Pensa faire tout bas le souhait d'être Veuf, \\
Et peut-être entre nous ne pouvoit-il mieux faire. \\
Les hommes, disoit-il, pour souffrir sont bien nez. \\
Peste soit du Boudin, \& du Boudin encore. \\
\vin Plût à Dieu, maudite Pecore, \\
\vin Qu'il te pendît au bout du nez!

La Priere aussitôt du Ciel fut écoutée, \\
Et dés que le Mari la parole lâcha \\
\vin Au nez de l'Epouse irritée \\
\vin L'Aune de Boudin s'attacha. \\
Ce prodige impréveu grandement le fâcha. \\
La Femme étoit jolie, elle avoit bonne grace, \\
Et pour dire sans fard la verité du fait, \\
\vin Cet ornement en cette place \\
\vin Ne faisoit pas un bon effet, \\
Si ce n'est qu'en pendant sur le bas du visage \\
Et lui fermant la bouche à tout moment \\
\vin Il l'empéchoit de parler aisément, \\
\vin Pour un Epoux merveilleux avantage. \\
Je pourrois bien, disoit-il à part soi \\
Pour me dédommager d'un malheur si funeste, \\
\vin Avec le souhait qui me reste \\
\vin Tout d'un plein saut me faire Roi, \\
Rien n'égale, il est vrai, la grandeur Souveraine, \\
\vin Mais encore faut-il songer \\
\vin Comment seroit faite la Reine, \\
Et dans quelle douleur ce seroit la plonger, \\
\vin De l'aller placer sur un Trone \\
\vin Avec un nez plus long qu'une aune. \\
\vin Il faut l'écouter sur cela; \\
\vin Et qu'elle même elle soit la Maîtresse \\
\vin De devenir une grande Princesse, \\
\vin En conservant l'horrible nez qu'elle a, \\
\vin Ou de demeurer Bucheronne, \\
\vin Avec un nez comme une autre personne, \\
Et tel qu'elle l'avoit avant ce malheur-là.

\vin La chose bien examinée, \\
Quoi qu'elle sçeût d'un Sceptre \& le prix \& l'effet, \\
\vin Et que quand on est couronnée \\
\vin On a toujours le nez bien fait, \\
Comme au desir de plaire il n'est rien qui ne cede, \\
\vin Elle aima mieux garder son Bavolet, \\
\vin Que d'être Reine \& d'être laide. \\
Ainsi le Bucheron ne changea point d'état; \\
\vin Il ne devint point Potentat, \\
\vin D'écus il n'emplit point sa Bourse, \\
Trop heureux d'employer le souhait qui restoit, \\
\vin Fraîle bonheur, pauvre ressource, \\
A remettre sa Femme en l'état qu'elle étoit; \\
\vin Tant il est vrai qu'aux hommes miserables, \\
Aveugles, imprudens, inquiéts, variables, \\
\vin Pas n'appartient de faire des souhaits, \\
\vin Et que peu d'entre eux sont capables \\
De bien user des dons que le Ciel leur a faits.
\end{verse}

% === POEMA 43 ===
\begin{verse}
En vous offrant, jeune \& sage Beauté \\
\vin Ce modele de patience, \\
\vin Je ne me suis jamais flatté \\
Que par vous de tout point il seroit imité \\
\vin C'en seroit trop en conscience.

\vin Mais Paris où l'homme est poli, \\
\vin Où le beau sexe né pour plaire \\
\vin Trouve son bonheur accompli, \\
\vin De tous côtez est si rempli \\
\vin D'Exemples du vice contraire, \\
\vin Qu'on ne peut en toute saison \\
\vin Pour s'en garder ou s'en défaire, \\
\vin Avoir trop de contrepoison.

\vin Une Dame aussi patiente \\
Que celle dont ici je reléve le prix, \\
\vin Seroit par tout une chose étonnante. \\
\vin Mais ce seroit un prodige à Paris. \\
\vin Les femmes y sont souveraines, \\
\vin Tout s'y regle selon leurs vœux, \\
\vin Enfin c'est un climat heureux \\
\vin Qui n'est habité que de Reines.

\vin Ainsi je voi que de toutes façons, \\
\vin Griselidis y sera peu prisée, \\
Et qu'elle y donnera matiere de risée, \\
\vin Par ses trop antiques leçons.

\vin Ce n'est pas que la patience \\
Ne soit une vertu des Dames de Paris, \\
Mais, par un long usage elles ont la science \\
De la faire exercer par leurs propres Maris.
\end{verse}

% === POEMA 44 ===
\begin{verse}
Au pié des celebres Montagnes \\
\vin Où le Pô s'échappant de dessous ses roseaux, \\
\vin Va dans le sein des prochaines Campagnes, \\
\vin Promener ses naissantes eaux, \\
\vin Vivoit un jeune et vaillant Prince, \\
\vin Les delices de sa Province. \\
Le Ciel en le formant, sur lui tout à la fois, \\
\vin Versa ce qu'il a de plus rare, \\
Ce qu'entre ses Amis d'ordinaire il separe, \\
\vin Et qu'il ne donne qu'aux grands Rois.

Comblé de tous les dons \& du corps \& de l'Ame, \\
Il fut robuste, adroit, propre au métier de Mars, \\
Et par l'instinct secret d'une divine flâme, \\
\vin Avec ardeur il aima les beaux arts. \\
Il aima les combats, il aima la Victoire, \\
\vin Les grands projets, les actes Valeureux, \\
Et tout ce qui fait vivre un beau nom dans l'Histoire; \\
\vin Mais son cœur tendre \& genereux \\
Fut encor plus sensible à la solide gloire \\
\vin De rendre ses peuples heureux. \\
\vin Ce temperament Héroïque \\
\vin Fut obscurci d'une sombre vapeur \\
\vin Qui chagrine \& melancolique, \\
\vin Lui faisoit voir dans le fond de son Cœur, \\
\vin Tout le beau sexe infidelle \& trompeur. \\
Dans la femme, où brilloit le plus rare merite, \\
\vin Il voyoit une ame hipocrite, \\
\vin Un Esprit d'orgueïl enivré, \\
Un cruel ennemi qui sans cesse n'aspire \\
\vin Qu'à prendre un souverain Empire \\
Sur l'Homme malheureux qui lui sera livré.

\vin Le frequent usage du Monde, \\
Où l'on ne voit qu'Epoux subjuguez ou trahis, \\
\vin Joint à l'air jaloux du Païs, \\
\vin Accrut encor cette haine profonde. \\
\vin Il jura donc plus d'une fois \\
Que quand même le Ciel pour lui plein de tendresse, \\
\vin Formeroit une autre Lucrece, \\
Jamais de l'himenée il ne suivroit les Loix. \\
Ainsi, quand le matin, qu'il donnoit aux affaires, \\
\vin Il avoit reglé sagement \\
\vin Toutes les choses necessaires \\
\vin Au bonheur du Gouvernement, \\
Que du foible orphelin, de la veuve oppressée, \\
\vin Il avoit conservé les droits, \\
Ou banni quelque impôt qu'une guerre forcée \\
\vin Avoit introduit autrefois; \\
\vin L'autre moitié de la journée \\
\vin A la Chasse étoit destinée, \\
\vin Ou les Sangliers \& les Ours, \\
\vin Malgré leur fureur \& leurs Armes \\
\vin Lui donnoient encor moins d'allarmes \\
Que le sexe charmant qu'il évitoit toujours.

Cependant ses sujets que leur interét presse \\
\vin De s'asseurer d'un Successeur \\
Qui les gouverne un jour avec même douceur, \\
A leur donner un fils le convioient sans cesse.

Un jour dans le Palais ils vinrent tous en corps \\
\vin Pour faire leurs derniers efforts; \\
\vin Un Orateur d'une grave apparence, \\
\vin Et le meilleur qui fût alors, \\
Dit tout ce qu'on peut dire en pareille occurrence \\
\vin Il marqua leur desir pressant \\
De voir sortir du Prince une heureuse Lignée \\
Qui rendit à jamais leur Etat florissant, \\
\vin Il lui dit même en finissant \\
\vin Qu'il voyoit un astre naissant \\
\vin Issu de son chaste hymenée \\
\vin Qui faisoit pâlir le croissant.

D'un ton plus simple \& d'une voix moins forte \\
Le Prince à ses sujets repondit de la sorte.

\vin Le zele ardent, dont je voi qu'en ce jour \\
\vin Vous me portez aux nœuds du mariage, \\
\vin Me fait plaisir, \& m'est de vôtre Amour \\
\vin Un agreable témoignage; \\
\vin J'en suis sensiblement touché, \\
Et voudrais dés demain pouvoir vous satisfaire \\
\vin Mais a mon sens l'Hymen est une affaire \\
Où plus l'homme est prudent, plus il est empêché. \\
\vin Observez bien toutes les jeunes filles; \\
\vin Tant qu'elles sont au sein de leurs familles \\
\vin Ce n'est que vertu, que bonté, \\
\vin Que pudeur, que sincerité; \\
\vin Mais sitôt que le mariage \\
\vin Au deguisement a mis fin, \\
\vin Et qu'ayant fixé leur destin \\
\vin Il n'importe plus d'être sage, \\
\vin Elles quittent leur personnage, \\
\vin Non sans avoir beaucoup pati, \\
\vin Et chacune dans son ménage \\
\vin Selon son gré prend son parti.

L'une d'humeur chagrine, \& que rien ne recrée, \\
\vin Devient une devote outrée, \\
\vin Qui crie \& gronde à tous momens, \\
\vin L'autre se façonne en Coquette, \\
\vin Qui sans cesse écoute ou caquette, \\
\vin Et n'a jamais assez d'Amans; \\
Celle ci des beaux arts follement curieuse, \\
\vin De tout décide avec hauteur, \\
\vin Et critiquant le plus habile autheur, \\
\vin Prend la forme de Precieuse; \\
\vin Cette autre s'erige en joüeuse, \\
Perd tout, argent, bijoux, bagues, meubles de prix, \\
\vin Et même jusqu'à ses habits. \\
Dans la diversité des routes qu'elles tiennent \\
\vin Il n'est qu'une chose où je voi \\
\vin Qu'enfin toutes elles conviennent, \\
\vin C'est de vouloir donner la Loi.

Or je suis convaincu que dans le mariage \\
\vin On ne peut jamais vivre heureux, \\
\vin Quand on y commande tous deux. \\
Si donc vous souhaittez qu'à l'Himen je m'engage, \\
\vin Cherchez une jeune Beauté \\
\vin Sans orgueil \& sans vanité, \\
\vin D'une obeïssance achevée, \\
\vin D'une patience éprouvée, \\
\vin Et qui n'ait point de volonté, \\
\vin Je la prendrai quand vous l'aurez trouvée.

Le prince, ayant mis fin à ce discours moral, \\
\vin Monte brusquement à cheval, \\
\vin Et court joindre à perte d'haleine \\
Sa meutte qui l'attend au milieu de la plaine.

Aprés avoir passé des prés \& des guerets, \\
Il trouve ses chasseurs couchez sur l'herbe verte \\
\vin Tous se levent, \& tous alerte, \\
Font trembler de leurs cors les hôtes des forêts. \\
\vin Des chiens courans, l'abboyante famille, \\
\vin Deçà, de là, parmi le chaume brille, \\
\vin Et les Limiers à l'œil ardent \\
Qui du fort de la bête à leur poste reviennent, \\
\vin Entraînent en les regardant \\
\vin Les forts valets qui les retiennent.

\vin S'étant instruit par un des siens \\
\vin Si tout est prêt, si l'on est sur la trace \\
Il ordonne aussitôt qu'on commence la chasse, \\
\vin Et fait donner le Cerf aux chiens. \\
\vin Le son des cors qui retentissent, \\
\vin Le bruit des chevaux qui hennissent \\
Et des chiens animez les pénétrans abois, \\
Remplissent la fôret de tumulte \& de trouble, \\
Et pendant que l'echo sans cesse les redouble, \\
S'enfonçent avec eux dans les plus creux du bois.

Le Prince par hasard ou par sa destinée, \\
\vin Prit une route détournée \\
\vin Où nul des chasseurs ne le suit; \\
\vin Plus il court, plus il s'en sépare: \\
\vin Enfin, à tel point il s'egare, \\
Que des chiens \& des cors il n'entend plus le bruit.

L'Endroit où le mena sa bijarre avanture, \\
\vin Clair de ruisseaux \& sombre de verdure, \\
Saisissoit les Esprits d'une secrette horreur; \\
\vin La simple \& naïve nature \\
\vin S'y faisoit voir \& si belle \& si pure, \\
\vin Que mille fois il benit son erreur.

\vin Rempli des douces rêveries \\
Qu'inspirent les grands bois, les eaux \& les prairies, \\
Il sent soudain frapper \& son cœur \& ses yeux \\
\vin Par l'objet le plus agreable, \\
\vin Le plus doux \& le plus aimable \\
\vin Qu'il eut jamais vu sous les Cieux. \\
\vin C'étoit une jeune Bergere \\
\vin Qui filoit aux bords d'un ruisseau, \\
\vin Et qui conduisant son troupeau, \\
\vin D'une main sage \& menagere \\
\vin Tournoit son agile fuzeau. \\
Elle auroit pû dompter les cœurs les plus sauvages; \\
\vin Des Lys, son teint a la blancheur, \\
\vin Et sa naturelle fraîcheur \\
S'étoit toûjours sauvée à l'ombre des boccages: \\
Sa bouche, de l'enfance avoit tout l'agrément, \\
Et ses yeux qu'adoucit une brune paupiere, \\
\vin Plus bleus que n'est le firmament, \\
\vin Avoient aussi plus de lumiere.

Le Prince, avec transport, dans le bois se glissant, \\
Contemple les beautez dont son Ame est émeüe, \\
\vin Mais le bruit qu'il fait en passant \\
De la belle sur lui fit détourner la veüe; \\
\vin Des qu'elle se vit apperçüe, \\
D'un brillant incarnat la prompte \& vive ardeur, \\
\vin De son beau teint redoubla la splendeur, \\
\vin Et sur son visage épandeüe, \\
\vin Y fit triompher la pudeur.

Sous le voile innocent de cette honte aimable, \\
Le Prince découvrit une simplicité, \\
\vin Une douceur, une sincerité, \\
\vin Dont il croyoit le beau sexe incapable, \\
\vin Et qu'il voyait dans toute leur beauté.

Saisi d'une frayeur pour lui toute nouvelle, \\
Il s'approche interdit, \& plus timide qu'elle, \\
\vin Lui dit d'une tremblante voix, \\
Que de tous ses veneurs il a perdu la trace, \\
\vin Et lui demande si la chasse \\
\vin N'a point passé quelque part dans le bois. \\
Rien n'a paru, Seigneur, dans cette solitude, \\
Dit-elle, \& nul ici que vous seul n'est venu; \\
\vin Mais n'ayez point d'inquiétude, \\
Je remettrai vos pas sur un chemin connu.

\vin De mon heureuse destinée \\
Je ne puis, lui dit-il, trop rendre grace aux Dieux, \\
\vin Depuis long-tems je frequente ces lieux, \\
Mais j'avois ignoré jusqu'à cette journée \\
\vin Ce qu'ils ont de plus precieux.

Dans ce tems elle voit que le Prince se baisse \\
\vin Sur le moitte bord du ruisseau, \\
\vin Pour étancher dans le cours de son eau \\
\vin La soif ardente qui le presse; \\
\vin Seigneur, attendez un moment, \\
\vin Dit-elle, \& courant promptement \\
\vin Vers sa cabane, elle y prend une tasse, \\
\vin Qu'avec joye \& de bonne grace, \\
\vin Elle presente à ce nouvel Amant.

Les vases precieux de cristal \& d'agathe \\
\vin Où l'or en mille endroits éclatte, \\
Et qu'un art curieux avec soin façonna: \\
N'eurent jamais pour lui, dans leur pompe inutile, \\
\vin Tant de beauté que le vase d'argile \\
\vin Que la Bergere lui donna.

Cependant pour trouver une route facile, \\
\vin Qui mene le Prince à la Ville, \\
Ils traversent des bois, des rochers escarpez \\
\vin Et de torrents entrecoupez, \\
Le Prince n'entre point dans de route nouvelle \\
Sans en bien observer, tous les lieux d'alentour; \\
\vin Et son ingénieux Amour \\
\vin\vin Qui songeoit au retour \\
\vin En fit une carte fidelle. \\
\vin Dans un boccage sombre \& frais \\
\vin Enfin la Bergere le meine, \\
\vin Où, de dessous ses branchages épais \\
\vin Il voit au loin dans le sein de la plaine \\
\vin Les toits dorez de son riche Palais.

\vin S'étant separé de la Belle, \\
\vin Touché d'une vive douleur, \\
\vin A pas lents il s'éloigne d'elle \\
\vin Chargé du trait qui lui perce le cœur. \\
\vin Le souvenir de sa tendre avanture, \\
\vin Avec plaisir le conduisit chez lui, \\
Mais dés le lendemain il sentit sa blessure, \\
Et se vit accablé de tristesse \& d'ennui.

\vin Dés qu'il le peut il retourne à la chasse, \\
\vin Où de sa suite adroitement \\
\vin Il s'échappe \& se débarrasse \\
\vin Pour s'égarer heureusement. \\
Des arbres \& des monts les cimes élevées. \\
\vin Qu'avec grand soin il avoit observées, \\
Et les avis secrets de son fidelle amour, \\
Le guiderent si bien que malgré les traverses, \\
\vin\vin De cent routes diverses, \\
De sa jeune Bergere il trouva le séjour. \\
Il sçut qu'elle n'a plus que son pere avec elle, \\
\vin Que Griselidis on l'appelle, \\
Qu'ils vivent doucement du lait de leurs brebis, \\
Et que de leur toison qu'elle seule elle file, \\
\vin Sans avoir recours à la Ville, \\
\vin Ils font eux-mêmes leurs habits.

\vin Plus il la voit plus il s'enflâme \\
\vin Des vives beautez de son ame. \\
Il connoit en voyant tant de dons précieux, \\
\vin Que si sa Bergere est si belle, \\
\vin C'est qu'une legere étincelle, \\
De l'esprit qui l'anime a passé dans ses yeux.

\vin Il ressent une joye extréme, \\
D'avoir si bien placé ses premieres amours, \\
Ainsi sans plus tarder, il fit dés le jour même \\
Assembler son Conseil \& lui tint ce discours.

\vin Enfin aux Loix de l'hyménée \\
\vin Suivant vos vœux je me vais engager, \\
Je ne prens point ma femme en païs étranger. \\
Je la prends parmi vous, belle, sage, bien née, \\
Ainsi que mes ayeux  ont fait plus d'une fois, \\
\vin Mais j'attendrai cette grande journée \\
\vin A vous informer de mon choix.

\vin Dés que la nouvelle fut sçüe, \\
\vin Partout elle fut répanduë. \\
\vin On ne peut dire avec combien d'ardeur \\
\vin L'allegresse publique \\
\vin De tous côtez s'explique; \\
\vin Le plus content fût l'Orateur, \\
\vin Qui par son discours pathetique \\
Croyoit d'un si grand bien être l'unique Auteur, \\
\vin Qu'il se trouvoit homme de consequence! \\
Rien ne peut resister à la grande éloquence, \\
\vin Disoit-il sans cesse en son cœur.

Le plaisir fut de voir le travail inutile, \\
\vin Des Belles de toute la Ville \\
\vin Pour s'attirer \& mériter le choix \\
Du Prince leur Seigneur, q'un air chaste \& modeste, \\
Charmoit uniquement \& plus que tout le reste, \\
\vin Ainsi qu'il l'avoit dit cent fois.

D'habit \& de maintien toutes elles changerent, \\
\vin D'un ton devot elles tousserent, \\
\vin Elles radoucirent leurs voix, \\
\vin De demi pied les coëffures baisserent, \\
La gorge se couvrit, les manches s'allongerent, \\
A peine on leur voyoit le petit bout des doigts.

\vin Dans la Ville avec diligence, \\
\vin Pour l'hymen dont le jour s'avance, \\
\vin On voit travailler tous les arts, \\
\vin Ici se font de magnifiques chars \\
\vin D'une forme toute nouvelle, \\
\vin Si beaux \& si bien inventez, \\
\vin Que l'or qui par tout étincelle, \\
\vin En fait la moindre des beautez.

Là, pour voir aisément \& sans aucun obstacle, \\
\vin Toute la pompe du spectacle, \\
\vin On dresse de longs échaffaux, \\
\vin Ici de grands Arcs triomphaux, \\
Où du Prince guerrier se celebre la gloire, \\
Et de l'amour sur lui l'éclatante victoire.

\vin Là sont forgez d'un art industrieux, \\
Ces feux qui par les coups d'un innocent Tonnerre, \\
\vin En effrayant la Terre, \\
De mille astres nouveaux embellissent les Cieux.

\vin Là d'un ballet ingenieux \\
Se concerte avec soin l'agreable folie, \\
Et là d'un Opéra peuplé de mille dieux, \\
Le plus beau que jamais ait produit l'Italie, \\
On entend repeter les Airs melodieux.

\vin Enfin, du fameux hymené, \\
\vin Arriva la grande journée.

\vin Sur le fond d'un Ciel vif \& pur, \\
\vin A peine l'Aurore vermeille, \\
\vin Confondoit l'or avec l'azur, \\
Que par tout en sursaut le beau sexe s'eveille; \\
Le peuple curieux s'épand de tous côtez, \\
En differens endroits des Guardes sont postez, \\
\vin Pour contenir la populace, \\
\vin Et la contraindre à faire place. \\
\vin Tout le Palais retentit de clairons, \\
De flutes, de hautbois, de rustiques musettes, \\
\vin Et l'on n'entend aux environs \\
\vin Que des tambours \& des trompettes. \\
Enfin le Prince sort entouré de sa Cour, \\
\vin Il s'éleve un long cri de joye, \\
Mais on est bien surpris quand au premier détour, \\
De la forêt prochaine on voit qu'il prend la voye, \\
\vin Ainsi qu'il faisoit chaque jour. \\
\vin Voilà, dit-on, son penchant qui l'emporte, \\
Et de ses passions, en dépit de l'amour, \\
\vin La Chasse est toûjours la plus forte.

\vin Il traverse rapidement \\
Les guerets de la plaine, \& gagnant la montagne, \\
Il entre dans le bois au grand étonnement \\
\vin De la Troupe qui l'accompagne. \\
Après avoir passé par différens détours, \\
Que son cœur amoureux se plaît à reconnaître, \\
\vin Il trouve enfin la cabane champêtre \\
\vin Où logent ses tendres amours.

\vin Griselidis de l'hymen informée, \\
\vin Par la voix de la Renommée, \\
\vin En avoit pris son bel habillement; \\
Et pour en aller voir la pompe magnifique \\
\vin De dessous sa case rustique \\
\vin Sortoit en ce même moment.

\vin Où courez-vous, si prompte \& si legere? \\
\vin Lui dit le prince en l'abordant, \\
Cessez de vous hâter, trop aimable Bergere, \\
La Nopce où vous allez, \& dont je suis l'Epoux, \\
\vin Ne saurait se faire sans vous.

\vin Oüi, je vous aime, \& je vous ai choisie \\
\vin Entre mille jeunes beautez \\
Pour passer avec vous le reste de ma vie. \\
Si toutefois mes vœux ne sont pas rejettez.

Ah! dit-elle, Seigneur, je n'ai garde de croire \\
Que je sois destinée à ce comble de gloire, \\
\vin Vous cherchez à vous divertir. \\
\vin Non, non, dit-il, je suis sincere, \\
\vin J'ai deja pour moi vôtre Pere. \\
(Le Prince avoit eu soin de l'en faire avertir) \\
\vin Daignez Bergere y consentir, \\
\vin C'est-là tout ce qui reste à faire. \\
Mais afin qu'entre nous une solide paix \\
\vin Eternellement se maintienne, \\
Il faudroit me jurer que vous n'aurez jamais \\
\vin D'autre volonté que la mienne.

Je le jure, dit-elle, \& je vous le promets; \\
Si j'avois épouzé le moindre du Village, \\
\vin J'obeïrois, son joug me serait doux, \\
\vin Hélas! combien donc davantage, \\
\vin Si je viens à trouver en vous, \\
\vin Et mon Seigneur et mon Epoux.

\vin Ainsi le Prince se déclare, \\
Et pendant que la Cour applaudit à son choix, \\
Il porte la Bergere à souffrir qu'on la pare \\
Des ornemens qu'on donne aux Epouzes des Rois. \\
Celles qu'à cet emploi leur devoir interesse, \\
Entrent dans la Cabane, \& là diligemment \\
Mettent tout leur savoir \& toute leur adresse \\
\vin A donner de la grace à chaque ajustement. \\
\vin Dans cette hutte où l'on se presse, \\
\vin Les Dames admirent sans cesse \\
\vin Avec quel art la pauvreté \\
\vin S'y cache sous la propreté; \\
\vin Et cette rustique Cabane, \\
Que couvre \& refraichit un spacieux Platane, \\
\vin Leur semble un séjour enchanté.

Enfin, de ce Reduit sort pompeuse \& brillante \\
\vin La Bergere Charmante, \\
\vin Ce ne sont qu'applaudissemens \\
\vin Sur sa beauté, sur ses habillemens; \\
\vin Mais sous cette pompe étrangere, \\
Déja plus d'une fois le Prince a regretté \\
\vin Des ornemens de la Bergere \\
\vin L'innocente simplicité.

\vin Sur un grand char d'or \& d'Ivoire \\
La Bergere s'assied pleine de Majesté, \\
\vin Le Prince y monte avec fierté, \\
\vin Et ne trouve pas moins de gloire \\
A se voir comme Amant assis à son côté, \\
Qu'à marcher en triomphe aprés une victoire; \\
\vin La Cour les suit \& tous gardent le rang \\
Que leur donne leur charge ou l'éclat de leur sang.

La Ville dans les champs presque toute sortie \\
\vin Couvroit les plaines d'alentour, \\
\vin Et du choix du Prince avertie, \\
Avec impatience attendoit son retour, \\
Il paroit, on le joint. Parmi l'épaisse foule \\
Du peuple qui se fend le char à peine roule; \\
Par les longs cris de joye à tout coup redoublez, \\
\vin Les chevaux émûs et troublez, \\
\vin Se cabrent, trépignent, s'élancent \\
\vin Et reculent plus qu'ils n'avancent.

\vin Dans le Temple on arrive enfin, \\
\vin Et là par la chaîne éternelle \\
\vin D'une promesse solennelle, \\
\vin Les deux Epoux unissent leur destin: \\
\vin Ensuite au Palais ils se rendent, \\
\vin Où mille plaisirs les attendent, \\
Où la Danse, les Jeux, les Courses, les Tournois \\
Repandent l'allegresse en differens endroits; \\
\vin Sur le soir le blond hymenée, \\
De ses chastes douceurs couronna la journée.

\vin Le lendemain les differents Etats \\
\vin\vin De toute la Province \\
Accourent haranguer la Princesse \& le Prince \\
\vin Par la voix de leurs Magistrats.

\vin De ses Dames environnée, \\
\vin Griselidis, sans paroître étonnée, \\
\vin En Princesse les entendit, \\
\vin En Princesse leur répondit. \\
Elle fit toute chose avec tant de prudence, \\
Qu'il sembla que le Ciel eût versé ses thrésors, \\
\vin Avec encor plus d'abondance \\
\vin Sur son Ame que sur son corps. \\
\vin Par son Esprit, par ses vives lumières, \\
Du Grand monde aussitôt elle prit les maniéres, \\
\vin Et même dés le premier jour \\
Des talens, de l'humeur des Dames de la Cour, \\
\vin Elle se fit si bien instruire, \\
\vin Que son bon sens jamais embarrassé \\
\vin Eut moins de peine à les conduire, \\
\vin Que ses brebis du tems passé.

Avant la fin de l'an des fruits de l'hymenée, \\
\vin Le Ciel benit leur couche fortunée, \\
Ce ne fut point un Prince, on l'eût bien souhaitté; \\
Mais la jeune Princesse avoit tant de beauté, \\
Que l'on ne songea plus qu'à conserver sa vie; \\
Le Pere qui lui trouve un air doux \& charmant, \\
\vin La venoit voir de moment en moment, \\
\vin Et la Mere encor plus ravie \\
\vin La regardoit incessamment.

\vin Elle voulut la nourrir elle-même, \\
Ah! dit-elle, comment m'exempter de l'emploi \\
\vin Que ses cris demandent de moi, \\
\vin Sans une ingratitude extrême; \\
\vin Par un motif de Nature ennemi \\
Pourrois-je bien vouloir de mon Enfant que j'aime, \\
\vin N'être la Mere qu'à demi. \\
Soit que le Prince eût l'ame un peu moins enflammé \\
\vin Qu'aux premiers jours de son ardeur, \\
\vin Soit que de sa maligne humeur \\
\vin La masse se fût rallumée, \\
\vin Et de son épaisse fumée \\
Eût obscurci ses sens et corrompu son Cœur; \\
\vin Dans tout ce que fait la Princesse, \\
Il s'imagine voir peu de sincerité, \\
\vin Sa trop grande vertu le blesse, \\
C'est un piege qu'on rend à sa credulité; \\
Son Esprit inquiet \& de trouble agité \\
\vin Croit tous les soupçons qu'il écoute, \\
\vin Et prend plaisir à revoquer en doute \\
\vin L'excez de sa felicité.

Pour guerir les chagrins dont son ame est atteinte \\
Il la suit, il l'observe, il aime à la troubler \\
\vin Par les ennuys de la contrainte, \\
\vin Par les alarmes de la crainte, \\
\vin Par tout ce qui peut deméler \\
\vin La verité d'avec la feinte. \\
\vin C'est trop, dit-il, me laisser endormir, \\
\vin Si ses vertus sont veritables \\
\vin Les traitemens les plus insupportables, \\
\vin Ne feront que les affermir.

\vin Dans son Palais il la tient reserrée, \\
Loin de tous les plaisirs qui naissent à la Cour, \\
Et dans sa chambre, où seule elle vit retirée, \\
\vin A peine il laisse entrer le jour. \\
\vin Persuadé que la Parure \\
\vin Et le superbe ajustement \\
Du sexe, que pour plaire a formé la Nature \\
\vin Est le plus doux enchantement. \\
\vin Il lui demande avec rudesse \\
Les Perles, les Rubis, les Bagues, les Bijoux \\
\vin Qu'il lui donna pour marque de tendresse, \\
Lorsque de son Amant il devint son Epoux.

\vin Elle dont la vie est sans tache, \\
\vin Et qui n'a jamais eu d'attache \\
\vin Qu'à s'acquiter de son devoir, \\
\vin Les lui donne sans s'émouvoir. \\
Et même le voyant se plaire à les reprendre, \\
\vin N'a pas moins de joye à les rendre \\
\vin Qu'elle en eût à les recevoir.

\vin Pour m'éprouver mon Epoux me tourmente, \\
Dit-elle, \& je voi bien qu'il ne me fait souffrir, \\
Qu'afin de reveiller ma vertu languissante, \\
Qu'un doux \& long repos pourrait faire perir. \\
S'il n'a pas ce dessein, du moins suis-je assurée \\
Que telle est du Seigneur la conduite sur moi; \\
Et que de tant de maux l'ennuyeuse durée, \\
N'est que pour exercer ma constance \& ma foi. \\
\vin Pendant que tant de malheureuses \\
\vin Errent au gré de leurs désirs; \\
\vin Par mille routes dangereuses. \\
\vin Aprés de faux \& vains plaisirs; \\
Pendant que le Seigneur dans sa lente Justice \\
\vin Les laisse aller au bord du précipice \\
\vin Sans prendre part à leur danger \\
Par un pur mouvement de sa bonté suprême \\
\vin Il me choisit comme un enfant qu'il aime \\
\vin Et s'applique à me corriger. \\
Aymons donc sa rigueur utilement cruelle \\
\vin On n'est heureux qu'autant qu'on a souffert; \\
\vin Aymons sa bonté paternelle \\
\vin Et la main dont elle se sert.

Le Prince a beau la voir obeïr sans contrainte \\
\vin A tous ses ordres absolus \\
Je voi le fondement de cette vertu feinte \\
Dit-il, \& ce qui rend tous mes coups superflus, \\
\vin C'est qu'ils n'ont porté leur atteinte \\
\vin Qu'à des endroits où son Amour n'est plus.

\vin Dans son Enfant, dans la jeune Princesse \\
\vin Elle a mis toute sa tendresse \\
\vin A l'éprouver si je veux reüssir \\
\vin C'est là qu'il faut que je m'adresse. \\
\vin C'est là que je puis m'éclaircir. \\
\vin Elle venait de donner la mamelle, \\
\vin Au tendre objet de son Amour ardent \\
Qui couché sur son sein se joüoit avec elle, \\
\vin Et rioit en la regardant:

Je voi que vous l'aymez, lui dit-il, cependant \\
Il faut que je vous l'ôte en cet âge encor tendre \\
Pour lui former les mœurs \& pour la preserver \\
De certains mauvais airs qu'avec vous l'on peut prendre; \\
\vin Mon heureux sort m'a fait trouver \\
Une Dame d'esprit qui saura l'élever \\
Dans toutes les vertus \& dans la politesse \\
\vin Que doit avoir une Princesse. \\
\vin Disposez-vous à la quitter \\
\vin On va venir pour l'emporter.

Il la laisse à ces mots, n'ayant pas le courage, \\
\vin Ni les yeux assez inhumains, \\
\vin Pour voir arracher de ses mains \\
\vin De leur Amour l'unique gage; \\
Elle de mille pleurs se baigne le visage, \\
\vin Et dans un morne accablement \\
Attend de son malheur le funeste moment.

Dés que d'une action si triste \& si cruelle \\
Le Ministre odieux à ses yeux se montra, \\
\vin Il faut obeïr lui dit-elle, \\
Puis prenant son Enfant qu'elle considera, \\
\vin Qu'elle baisa d'une ardeur maternelle, \\
Qui de ses petits bras tendrement la serra, \\
\vin Toute en pleurs elle le livra. \\
\vin Ah! que sa douleur fut amere! \\
\vin Arracher l'Enfant ou le Cœur \\
\vin Du sein d'une si tendre Mere, \\
\vin\vin C'est la même douleur.

\vin Prés de la Ville étoit un monastère, \\
\vin Fameux par son antiquité, \\
Où des vierges vivoient dans une regle austere, \\
Sous les yeux d'une Abbesse illustre en pieté. \\
\vin Ce fut là que dans le silence, \\
\vin Et sans déclarer sa naissance, \\
On déposa l'Enfant \& des bagues de prix, \\
\vin Sous l'espoir d'une recompense \\
\vin Digne de soins que l'on en auroit pris.

Le Prince qui tâchoit d'éloigner par la Chasse \\
\vin Le vif remords qui l'embarrasse \\
\vin Sur l'excez de sa cruauté, \\
\vin Craignoit de revoir la Princesse, \\
Comme on craint de revoir une fiere Tigresse \\
\vin A qui son faon vient d'être ôté: \\
\vin Cependant il en fut traité \\
\vin Avec douceur, avec caresse, \\
\vin Et même avec cette tendresse, \\
Qu'elle eut aux plus beaux jours de sa prosperité.

Par cette complaisance \& si grande \& si prompte, \\
\vin Il fut touché de regret \& de honte, \\
\vin Mais son chagrin demeura le plus fort: \\
Ainsi, deux jours aprés, avec des larmes feintes, \\
Pour lui porter encor de plus vives atteintes, \\
\vin Il lui vient dire que la mort \\
De leur aimable Enfant avoit fini le sort.

Le coup inopiné mortellement la blesse \\
\vin Cependant malgré sa tristesse, \\
Ayant veu son Epoux qui changeoit de couleur, \\
\vin Elle parut oublier son malheur, \\
\vin Et n'avoir même de tendresse \\
Que pour le consoler de sa fausse douleur.

\vin Cette bonté, cette ardeur sans égale \\
\vin D'amitié conjugale, \\
Du Prince tout à coup désarmant la rigueur \\
Le touche, le pénetre, \& lui change le Cœur, \\
\vin Jusques-là qu'il lui prend envie \\
\vin De déclarer que leur Enfant \\
\vin Joüit encore de la vie: \\
Mais sa bile s'éleve, \&, fiere lui defend \\
\vin De rien découvrir du mystere \\
\vin Qu'il peut-être utile de faire.

Dès ce bien heureux jour telle des deux Epoux \\
\vin Fut la mutuelle tendresse, \\
Qu'elle n'est point plus vive aux momens les plus doux \\
\vin Entre l'Amant \& la Maîtresse. \\
Quinze fois le soleil pour former les saisons, \\
Habita tour à tour dans ses douze maisons, \\
\vin Sans rien voir qui les desunisse: \\
\vin Que si quelques fois par caprice \\
\vin Il prend plaisir à la facher, \\
\vin C'est seulement pour empêcher \\
\vin Que l'amour ne se ralentisse, \\
Tel que le forgeron qui pressant son labeur \\
\vin Repand un peu d'eau sur la braize \\
\vin De sa languissante fournaise \\
\vin Pour en redoubler la chaleur.

\vin Cependant la jeune Princesse. \\
\vin Croissoit en esprit, en sagesse, \\
\vin A la douceur, à la naïveté \\
\vin Qu'elle tenoit de son aimable Mere, \\
\vin Elle joignit de son Illustre Pere \\
\vin L'agreable et noble fierté. \\
L'Amas de ce qui plaît dans chaque Caractère \\
\vin Fit une parfaite beauté.

\vin Par tout comme un Astre elle brille, \\
\vin Et par hazard un Seigneur de la Cour, \\
\vin Jeune, bien fait \& plus beau que le jour, \\
\vin L'ayant vû paroître à la Grille, \\
\vin Conçût pour elle un violent amour. \\
Par l'instinct qu'au beau sexe a donné la nature, \\
\vin Et que toutes les beautez ont, \\
\vin De voir l'invisible blessure \\
Que font leurs yeux, au moment qu'ils la font, \\
\vin La Princesse fut informée \\
\vin Qu'elle étoit tendrement aimée. \\
\vin Aprés avoir quelque tems resisté, \\
\vin Comme on le doit avant que de se rendre, \\
\vin D'un amour également tendre \\
\vin Elle l'aima de son côté.

\vin Dans cet Amant, rien n'étoit à reprendre \\
Il étoit beau, vaillant, né d'illustres Ayeux \\
\vin Et dés long-tems, pour en faire son Gendre \\
\vin Sur lui le Prince avoit jetté les yeux. \\
Ainsi donc avec joye il apprit la nouvelle, \\
\vin De l'ardeur tendre \& mutuelle \\
\vin Dont bruloient ces jeunes Amans, \\
\vin Mais il lui prit une bizarre envie, \\
De leur faire acheter par de cruels tourmens, \\
\vin Le plus grand bonheur de leur vie.

\vin Je me plairai, dit-il, à les rendre contens; \\
\vin Mais il faut que l'inquietude \\
\vin Par tout ce qu'elle a de plus rude, \\
\vin Rende encor leurs feux plus constans; \\
\vin De mon Epouze en même tems, \\
\vin J'exercerai la patience, \\
\vin Non point comme jusqu'à ce jour, \\
\vin Pour rasseurer ma folle défiance; \\
\vin Je ne dois plus douter de son amour: \\
Mais pour faire éclatter aux yeux de tout le monde, \\
Sa bonté, sa douceur, sa sagesse profonde; \\
Afin que de ses dons si grands, si precieux, \\
\vin La terre se voyant parée, \\
\vin En soit de respect pénétrée, \\
Et par reconnaissance en rende grace aux Cieux.

Il déclare en public que manquant de lignée, \\
En qui l'Etat un jour retrouve son Seigneur, \\
Que la fille qu'il eut de son fol hymenée \\
\vin Etant morte aussi-tôt que née, \\
\vin Il doit ailleurs chercher plus de bonheur. \\
Que l'Epouze qu'il prend est d'illustre naissance, \\
\vin Qu'en un Couvent on l'a jusqu'à ce jour \\
\vin Fait élever dans l'innocence, \\
Et qu'il va par l'hymen couronner son amour.

\vin On peut juger à quel point fut cruelle \\
Aux deux jeunes Amans cette affreuse nouvelle; \\
Ensuite sans marquer ni chagrin ni douleur, \\
\vin Il avertit son Epouze fidelle, \\
\vin Qu'il faut qu'il se separe d'elle \\
\vin Pour éviter un extreme malheur; \\
Que le peuple indigné de sa basse naissance \\
Le force à prendre ailleurs une digne alliance.

\vin Il faut, dit-il, vous retirer \\
\vin Sous vôtre toit de chaume \& de fougere \\
Aprés avoir repris vos habits de Bergere, \\
\vin Que je vous ai fait preparer. \\
Avec une tranquile \& muëtte constance, \\
La Princesse entendit prononcer sa sentence; \\
\vin Sous les dehors d'un visage serain \\
\vin Elle devoroit son chagrin, \\
Et sans que la douleur diminuât ses charmes, \\
\vin De ses beaux yeux tomboient de grosses larmes, \\
Ainsi que quelquefois au retour du Printems, \\
\vin Il fait soleil, \& pleut en même tems.

Vous êtes mon Epoux, mon Seigneur, \& mon Maître, \\
(Dit-elle en soûpirant, prête à s'évanoüir,) \\
Et quelque affreux que soit ce que je viens d'ouïr, \\
\vin Je saurai vous faire connoître \\
Que rien ne m'est si cher que de vous obeir.

Dans sa chambre aussi-tôt seule elle se retire \\
Et là se dépoüillant de ses riches habits, \\
\vin Elle reprend paisible \& sans rien dire, \\
\vin Pendant que son cœur en soûpire, \\
\vin Ceux qu'elle avoit en gardant ses Brebis. \\
\vin En cet humble \& simple équipage, \\
Elle aborde le Prince \& lui tient ce langage.

\vin Je ne puis m'éloigner de vous \\
\vin Sans le pardon d'avoir sû vous déplaire, \\
\vin Je puis souffrir le poids de ma misere, \\
Mais je ne puis, Seigneur, souffrir votre courroux. \\
Accordez cette grace à mon regret sincere, \\
Et je vivrai contente en mon triste séjour, \\
\vin Sans que jamais le tems altere \\
Ni mon humble respect, ni mon fidelle amour. \\
Tant de soumission, \& tant de grandeur d'ame \\
\vin Sous un si vil habillement, \\
Qui dans le cœur du Prince en ce même moment \\
Raveilla tous les traits de sa premiere flâme, \\
Alloient casser l'arrêt de son bannissement. \\
\vin Emû par de si puissants charmes, \\
\vin Et prêt à repandre des larmes, \\
\vin Il commençoit à s'avancer, \\
\vin\vin Pour l'embrasser.

\vin Quand tout à coup l'imperieuse gloire, \\
\vin D'être ferme en son sentiment \\
\vin Sur son amour remporta la victoire, \\
Et le fit en ces mots répondre durement.

De tout le temps passé j'ai perdu la mémoire, \\
\vin Je suis content de vôtre repentir, \\
\vin Allez il est tems de partir.

Elle part aussi-tôt, \& regardant son Pere, \\
Qu'on avoit revêtu de son rustique habit, \\
Et qui le cœur percé d'une douleur amere, \\
Pleuroit un changement si prompt \& si subit. \\
Retournons, lui dit-elle, en nos sombres boccages, \\
Retournons habiter nos demeures sauvages, \\
Et quittons sans regret la pompe des Palais, \\
Nos cabanes n'ont pas tant de magnificence, \\
\vin Mais on y voit régner dans l'innocence \\
Un plus ferme repos, une plus douce paix.

\vin Dans son desert a grand'peine arrivée, \\
\vin Elle reprend \& quenoüille \& fuzeaux, \\
\vin Et va filer au bord des mêmes eaux \\
\vin Où le Prince l'avoit trouvée. \\
\vin Là son cœur tranquille \& sans fiel, \\
\vin Cent fois le jour demande au Ciel, \\
Qu'il comble son Epoux de gloire, de richesses, \\
Et qu'à tous ses desirs il ne refuse rien. \\
\vin Un Amour nourri de caresses \\
\vin N'est pas plus ardent que le sien.

\vin Ce cher Epoux qu'elle regrette, \\
\vin Voulant encore l'éprouver, \\
\vin Lui fait dire dans sa retraite \\
\vin Qu'elle ait à le venir trouver.

Griselidis, dit-il, dés qu'elle se présente, \\
Il faut que la Princesse à qui je dois demain \\
\vin Dans le Temple donner la main, \\
\vin De vous \& de moi soit contente. \\
Je vous demande ici tous vos soins, \& je veux \\
Que vous m'aidiez à plaire à l'objet de mes vœux, \\
Vous savez de quel air il faut que l'on me serve, \\
\vin Point d'épargne, point de reserve, \\
Que tout sente le Prince, \& le Prince amoureux. \\
\vin Employez toute votre adresse \\
\vin A parer son appartement, \\
\vin Que l'abondance, la richesse, \\
\vin La propreté, la politesse \\
\vin S'y fasse voir également; \\
\vin Enfin songez incessamment \\
\vin Que c'est une jeune Princesse \\
\vin\vin Que j'aime tendrement.

\vin Pour vous faire entrer davantage \\
\vin Dans les soins de vôtre devoir, \\
\vin Je veux ici vous faire voir \\
Celle qu'à bien servir mon ordre vous engage. \\
\vin Telle qu'aux portes du Levant \\
\vin Se montre la naissante Aurore, \\
\vin Telle parut en arrivant \\
\vin La Princesse plus belle encore. \\
\vin Griselidis à son abord \\
Dans le fond de son cœur sentit un doux transport \\
\vin De la tendresse maternelle; \\
\vin Du tems passé, de ses jours bienheureux, \\
\vin Le souvenir en son cœur se rappelle, \\
\vin Helas! ma fille, en soi-même, dit-elle, \\
Si le Ciel favorable eût écouté mes vœux, \\
Seroit presqu'aussi grande \& peut-être aussi belle!

Pour la jeune Princesse en ce même moment, \\
Elle prit un amour si vif, si vehement, \\
\vin Qu'aussi-tôt qu'elle fut absente, \\
\vin En cette sorte au Prince elle parla, \\
Suivant sans le savoir, l'instinct qui s'en mêla. \\
\vin Souffrez, Seigneur, que je vous represente, \\
\vin Que cette Princesse charmante, \\
\vin Dont vous allez être l'Epoux, \\
Dans l'aise, dans l'éclat, dans la pourpre nourrie, \\
Ne pourra supporter, sans en perdre la vie, \\
Les mêmes traittements que j'ai reçû de vous. \\
\vin Le besoin, ma naissance obscure, \\
\vin M'avoient endurcie aux travaux \\
Et je pouvois souffrir toutes sortes de maux \\
\vin Sans peine \& même sans murmure; \\
Mais elle qui jamais n'a connu la douleur, \\
\vin Elle mourra dés la moindre rigueur, \\
Dés la moindre parole un peu seche, peu dure, \\
\vin Hélas! Seigneur, je vous conjure, \\
\vin De la traitter avec douceur.

Songez, lui dit le Prince avec un ton severe, \\
\vin A me servir selon votre pouvoir, \\
\vin Il ne faut pas qu'une simple Bergere \\
\vin Fasse des leçons, \& s'ingere, \\
\vin De m'avertir de mon devoir. \\
\vin Griselidis à ces mots sans rien dire, \\
\vin Baisse les yeux et se retire.

Cependant pour l'hymen les Seigneurs invitez, \\
\vin Arriverent de tous côtez, \\
\vin Dans une magnifique salle \\
\vin Où le Prince les assembla; \\
Avant que d'allumer la torche nuptiale, \\
\vin En cette sorte il leur parla. \\
\vin Rien au monde aprés l'esperance; \\
\vin N'est plus trompeur que l'apparence: \\
Ici l'on en peut voir un exemple éclatant, \\
\vin Qui ne croiroit que ma jeune Maîtresse, \\
\vin Que l'hymen va rendre Princesse, \\
Ne soit heureuse \& n'ait le cœur content? \\
\vin Il n'en est rien pourtant.

\vin Qui pourrait s'empêcher de croire, \\
Que ce jeune Guerrier amoureux de la gloire, \\
N'aime à voir cet hymen, lui qui dans les Tournois \\
Va sur tous ses Rivaux remporter la victoire, \\
\vin Cela n'est pas vrai toutefois.

Qui ne croiroit encor qu'en sa juste colere, \\
Griselidis ne pleure \& ne se desespere? \\
Elle ne se plaint point, elle consent à tout. \\
Et rien n'a pû pousser sa patience à bout.

Qui ne croiroit enfin que de ma destinée, \\
Rien ne peut égaler la course fortunée, \\
En voyant les appas de l'objet de mes vœux? \\
Cependant si l'hymen me lioit de ses nœuds, \\
\vin J'en concevrois une douleur profonde, \\
\vin Et de tous les Princes du monde, \\
\vin Je serois le plus malheureux.

L'énigme vous paroit difficile à comprendre, \\
\vin Deux mots vont vous la faire entendre, \\
\vin Et ces deux mots feront évanoüir \\
\vin Tous les malheurs que vous venez d'oüir.

Sachez, poursuivit-il, que l'aimable personne \\
\vin Que vous croyez m'avoir blessé le cœur, \\
\vin Est ma fille, \& que je la donne \\
\vin Pour femme à ce jeune Seigneur, \\
\vin Qui l'aime d'un amour extréme, \\
\vin Et dont il est aimé de même.

\vin Sachez encor, que touché vivement \\
\vin De la patience \& du zèle \\
\vin De l'Epouze sage \& fidelle \\
\vin Que j'ai chassée indignement, \\
\vin Je la reprens, afin que je repare, \\
Par tout ce que l'amour peut avoir de plus doux, \\
\vin Le traittement dur \& barbare \\
\vin Qu'elle a reçû de mon esprit jaloux. \\
\vin Plus grande sera mon étude, \\
\vin A prevenir tous ses desirs \\
\vin Qu'elle ne fut dans mon inquietude, \\
\vin A l'accabler de déplaisirs; \\
Et si dans tous les tems doit vivre la mémoire \\
Des ennuis dont son cœur ne fut point abattu, \\
Je veux que plus encore on parle de la gloire, \\
Dont j'aurai couronné sa supréme vertu.

\vin Comme quand un épais nuage \\
\vin\vin A le jour obscurci, \\
\vin Et que le Ciel de toutes parts noirci, \\
\vin Menace d'un affreux orage; \\
Si de ce voile obscur par les vents écarté, \\
\vin Un brillant rayon de clarté, \\
\vin Se repand sur le Païsage, \\
\vin Tout rit \& reprend sa beauté, \\
Telle dans tous les yeux où régnoit la tristesse \\
Eclatte tout à coup une vive allegresse.

\vin Par ce prompt éclaircissement \\
\vin La jeune Princesse ravie \\
D'apprendre que du Prince elle a reçû la vie, \\
Se jette à ses genoux qu'elle embrasse ardemment, \\
Son pere qu'attendrit une fille si chere, \\
La releve, la baise, \& la meine à sa mere, \\
A qui trop de plaisir en un même moment, \\
\vin Otoit presque tout sentiment. \\
\vin Son cœur qui tant de fois en proye \\
\vin Aux plus cuisans traits du malheur, \\
\vin Supporta si bien la douleur, \\
\vin Succombe au doux poids de la joye; \\
A peine de ses bras pouvoit-elle serrer \\
\vin L'aimable Enfant que le Ciel lui renvoye, \\
\vin Elle ne pouvoit que pleurer.

Assez dans d'autres tems vous pourrez satisfaire, \\
\vin Lui dit le Prince, aux tendresses du sang, \\
Reprenez les habits qu'exige votre rang, \\
\vin Nous avons des nopces à faire. \\
Au Temple on conduisit les deux jeunes Amans, \\
\vin Où la mutuelle promesse \\
\vin De se cherir avec tendresse, \\
Affermit pour jamais leurs doux engagemens, \\
Ce ne sont que plaisirs, que Tournois magnifiques, \\
\vin Que jeux, que dances, que musiques, \\
\vin Et que Festins delicieux, \\
Où sur Griselidis se tournent tous les yeux, \\
\vin Où sa patience éprouvée, \\
\vin Jusques au Ciel est élevée, \\
\vin Par mille éloges glorieux: \\
Des peuples réjoüis la complaisance est telle, \\
\vin Pour leur Prince capricieux; \\
Qu'ils vont jusqu'à loüer son épreuve cruelle, \\
\vin A qui d'une vertu si belle, \\
Si seante au beau sexe, \& si rare en tous lieux, \\
\vin On doit un si parfait modele.
\end{verse}

% === POEMA 45 ===
\begin{verse}
Est-ce une raison décisive \\
D'ôter un bon mets d'un repas, \\
Parce qu'il s'y trouve un convive \\
Qui par malheur ne l'aime pas? \\
Il faut que tout le monde vive, \\
Et que les mets, pour plaire à tous, \\
Soient differens comme les goûts.
\end{verse}
